\documentclass[11pt,a4paper]{article}

% ------------------------------------------------------------
% Packages
% ------------------------------------------------------------
\usepackage[margin=2.1cm]{geometry}
\usepackage{array}
\usepackage{booktabs}
\usepackage{longtable}
\usepackage{enumitem}
\usepackage{hyperref}
\usepackage[T1]{fontenc}
\usepackage{lmodern}
\usepackage{xcolor}
\usepackage{amsmath,amssymb}

% ------------------------------------------------------------
% Formatting
% ------------------------------------------------------------
\hypersetup{
    colorlinks=true,
    linkcolor=blue,
    urlcolor=blue,
    citecolor=blue
}

\setlength{\parindent}{0pt}
\setlength{\parskip}{0.6em}
\setlist[itemize]{leftmargin=*, topsep=0.2em, itemsep=0.15em}
\setlist[enumerate]{leftmargin=*, topsep=0.2em, itemsep=0.15em}

\newcommand{\weekheading}[2]{%
\section*{Week #1: #2}
\addcontentsline{toc}{section}{Week #1: #2}
}

\newcommand{\planningbox}[2]{%
\subsection*{#1}
#2
}

\newcommand{\phrase}[1]{%
\begin{quote}
\emph{#1}
\end{quote}
}

\newcommand{\stickingpoint}[1]{%
\par\smallskip
\noindent\textcolor{red!65!black}{\(\blacktriangleright\)\ \textbf{Likely sticking point.}} #1
\par\smallskip
}

\newcommand{\qualification}[1]{%
\par\smallskip
\noindent\textcolor{blue!70!black}{\(\dagger\)\ \textbf{Qualification.}} #1
\par\smallskip
}

\newcommand{\objectcheck}[1]{%
\par\smallskip
\noindent\textcolor{green!45!black}{\(\diamond\)\ \textbf{Object check.}} #1
\par\smallskip
}

% ------------------------------------------------------------
% Document
% ------------------------------------------------------------
\begin{document}

\begin{center}
    {\Large \textbf{Masterclass Narrative Planning Document}}\\[0.4em]
    {\large \textbf{Topological Data Analysis for Static and Dynamic Data}}\\[0.5em]
    \textit{A conceptual route from data to persistent homology}
\end{center}

\tableofcontents

\section*{Purpose and shape of the document}
\addcontentsline{toc}{section}{Purpose and shape of the document}

This document turns the week-by-week planning table into a more narrative
planning document. It is not yet a lecture script. Its purpose is to record the
story each week needs to tell, the mathematical objects that must be introduced,
the examples that should carry the intuition, and the phrases or explanatory
moves that are worth preserving.

The masterclass has two linked aims. The first is foundational: colleagues and
incoming Honours students should become well-versed enough in algebraic topology
and TDA to read, compute, question and adapt persistent homology pipelines. The
second is methodological: participants should be able to decide which
topological tool is appropriate for which data object and scientific question.
The Discovery Project motivates the masterclass, but reservoirs and swarms are
used as examples rather than as the curriculum itself.

\section*{Reading symbols}
\addcontentsline{toc}{section}{Reading symbols}

The following markers flag details that are easy to skate over but often become
the point at which an otherwise mathematically sophisticated audience loses the
thread.

\stickingpoint{Marks a conceptual transition or an apparently simple statement
that needs to be made concrete. These are the places most likely to prompt
``but why is that true?'' or ``what exactly is larger here?''}

\qualification{Marks a condition under which a statement is true, or a
distinction that should not be suppressed for convenience.}

\objectcheck{Stops to name the current objects and maps explicitly. In this
material, confusion often comes from moving silently between spaces, chain
spaces, homology groups, persistence modules and diagrams.}

\section*{Recurring questions}
\addcontentsline{toc}{section}{Recurring questions}

These questions should recur throughout the masterclass.

\begin{itemize}
    \item What is the mathematical object?
    \item What is the scientific object it is meant to represent?
    \item What information is being forgotten?
    \item What information is being added by the construction?
    \item What does birth mean under this filtration?
    \item What does death mean under this filtration?
    \item At what level are we working: space, homology, module, barcode or summary?
    \item Are the maps between objects part of the scientific information?
    \item Is feature identity required, or are summaries enough?
    \item Is there more than one genuine parameter?
    \item What null model or surrogate should be used?
    \item What simpler baseline must topology beat?
    \item This construction would be misleading if \ldots
\end{itemize}

\weekheading{1}{Shape, topology and the need for a computable invariant}

\planningbox{Story}{
The first week should make the case for topology as a lens without pretending
that topology is always the right lens. Begin with the distinction between a
directly given object and a constructed representation. A mug, donut, circle or
torus can be discussed as an object. A point cloud, scalar field, image, delay
embedding, interaction graph or reservoir-state cloud is a representation built
from observations, samples, relations or dynamics.

The central move is from metric geometry to topology. Geometry measures
distances, angles, curvature, scale and coordinates. Topology generalises the
idea of closeness by replacing exact distance with neighbourhood structure. In
a metric space, closeness is measured by a distance \(d(x,y)\). In a
topological space, closeness is encoded by open sets. Open sets are the regions
in which one can move a little without leaving the set.
}

\phrase{Openness in topology plays the role that distance plays in geometry:
it is the formal structure that says what closeness means.}

\planningbox{Mathematics to introduce}{
A topological space is a set \(X\) together with a collection \(\mathcal{T}\) of
subsets of \(X\), called open sets, such that
\begin{itemize}
    \item \(X\) and \(\varnothing\) are open;
    \item any union of open sets is open;
    \item any finite intersection of open sets is open.
\end{itemize}
The topology tells us what neighbourhoods, continuity, convergence,
connectedness and compactness mean. For this masterclass, the point is not to
develop point-set topology in detail, but to show that topology is a rule for
locality and continuity, not a vague synonym for shape.

A continuous function \(f:X\to \mathbb{R}\) assigns a real number to each point
of a topological space and changes without jumps, in the topological sense.
This language prepares the later use of sublevel sets.
}

\planningbox{Examples}{
Use the mug and donut to separate geometry from topology. Use the circle and
ellipse to show that topology preserves cyclic organisation while geometry
distinguishes metric form. Use the noisy circle versus noisy disk to introduce
the idea that homology asks a different question from precise geometry: does
the constructed representation organise around a one-dimensional hole?

The Datasaurus Dozen can motivate the general failure of scalar summaries, but
it should be used with a disclaimer. The point is not that topology is always
stronger than geometry. The point is that different summaries see different
structure.
}

\planningbox{DP relevance}{
For dynamical systems, the object of interest is often not directly visible. It
may be an attractor, recurrent orbit, collective regime, interaction structure,
state-space manifold or changing spatial pattern. The data are traces or
representations of that object. The DP question is whether topology can provide
computable witnesses of memory, separability, stability or collective
organisation. At this stage, the DP is only a reason to care about the lens.
}

\weekheading{2}{Simplicial complexes and homology}

\planningbox{Story}{
Week 2 moves from qualitative shape to a combinatorial object on which topology
can be computed. The central message is that a simplicial complex is the finite
bookkeeping device that lets topology become linear algebra. The complex is
treated as already chosen. The aim is to understand what homology computes once
\(K\) has been chosen.
}

\planningbox{Mathematics to introduce}{
A \(p\)-simplex is a \(p\)-dimensional building block: vertices, edges,
triangles, tetrahedra and their higher-dimensional analogues. A simplicial
complex \(K\) is a collection of simplices satisfying downward closure: if
\(\sigma\in K\) and \(\tau\subseteq\sigma\), then \(\tau\in K\). This condition
is not decoration. It ensures that whenever a simplex is present, all of its
faces are also present.

\objectcheck{At this stage, \(K\) is the topological/combinatorial object.
The vector spaces come next: one vector space \(C_p(K;\mathbb F)\) for the
\(p\)-simplices in each dimension.}

Over a field \(\mathbb{F}\), a \(p\)-chain is a formal linear combination of
\(p\)-simplices:
\[
c=\sum_i a_i\sigma_i,\qquad a_i\in\mathbb F.
\]
Thus \(C_p(K;\mathbb F)\) is a vector space whose basis elements are the
\(p\)-simplices. Over \(\mathbb{F}_2\), the coefficients are \(0\) or \(1\), so
a chain can be viewed as a subset of the \(p\)-simplices. Addition is symmetric
difference: simplices appearing twice cancel because \(1+1=0\) mod \(2\).

The boundary map lowers dimension and is linear:
\[
\partial_p:C_p(K;\mathbb F)\to C_{p-1}(K;\mathbb F).
\]
The boundary of an edge is its two endpoints; the boundary of a triangle is its
three edges; the boundary of a tetrahedron is its four triangular faces. The
identity \(\partial^2=0\) says that every boundary is automatically a cycle.

Cycles form the kernel
\[
Z_p=\ker \partial_p,
\]
and boundaries form the image
\[
B_p=\operatorname{im}\partial_{p+1}.
\]
Both are vector subspaces of \(C_p\), with \(B_p\subseteq Z_p\). Homology is
therefore the quotient vector space
\[
H_p(K;\mathbb F)=Z_p/B_p.
\]
Its elements are homology classes: cycles are identified when they differ by a
boundary. The Betti number is
\[
\beta_p=\dim H_p(K;\mathbb F).
\]

\stickingpoint{Homology is a vector space because the chain coefficients have
been chosen in a field, the boundary map is linear, and \(H_p=Z_p/B_p\) is a
quotient of vector spaces. A ``hole'' is represented by an equivalence class of
cycles, not by one privileged geometric loop.}

\qualification{With integer coefficients, \(H_p(K;\mathbb Z)\) is generally an
abelian group rather than a vector space. persistent homology computations
usually use field coefficients because this makes the algebra and interval
decomposition much cleaner.}
}

\phrase{Homology counts cycles that are not merely the boundary of something
higher-dimensional.}

\planningbox{Examples}{
Use an outline triangle and a filled triangle. In the outline triangle, the sum
of the three edges is a cycle with zero boundary. It is not the boundary of any
two-chain because no triangular face is present, so it represents a nonzero
class in \(H_1\). In the filled triangle, the same edge cycle is now
\(\partial_2(\text{triangle})\), so its homology class is zero.

Use a square loop as the simplest visible \(H_1\) class and a hollow tetrahedron
as a first \(H_2\) example. Ask explicitly which chains are cycles, which cycles
are boundaries, and which distinctions survive in the quotient.

For \(H_1(S^1)\), the winding number gives useful intuition. Loops around a
circle are classified by how many times they wind around the hole. Over
\(\mathbb{Z}\), this gives \(H_1(S^1)\cong \mathbb{Z}\). For the computational
pipeline, we usually work over a field such as \(\mathbb{F}_2\), but the
winding-number story helps explain why homology classes are algebraic
representatives of holes rather than the holes themselves.
}

\planningbox{DP relevance}{
Graphs enter as the simplest data-derived complexes. An interaction graph can
be treated as a one-dimensional simplicial complex, with vertices and edges
only. This supports component and graph-cycle homology. To study
higher-dimensional topology, the graph must be promoted to another object, such
as a clique complex. That promotion is a modelling assumption: pairwise mutual
interaction is being treated as evidence of a higher-order collective unit.
}

\weekheading{3}{Filtrations}

\planningbox{Story}{
Week 3 introduces the filtration. The key move is that one complex is not
enough for data analysis because a single threshold is usually arbitrary. A
filtration records how a complex grows as a meaningful parameter changes. The
parameter may be distance, density, intensity, similarity, interaction
strength, energy or time.
}

\planningbox{Mathematics to introduce}{
A filtration function assigns a real number to each simplex:
\[
f:K\to\mathbb{R}.
\]
It must be monotone on faces:
\[
\tau\subseteq\sigma \quad\Longrightarrow\quad f(\tau)\leq f(\sigma).
\]
This ensures that a face appears before any simplex that contains it. Sorting
simplices by filtration value gives a nested sequence
\[
\varnothing=K_0\subseteq K_1\subseteq\cdots\subseteq K_n=K,
\]
also denoted \(K_\bullet\).

\stickingpoint{Here ``\(K_j\) is larger than \(K_i\)'' means larger as a
simplicial complex: every simplex already in \(K_i\) is still present in
\(K_j\), and \(K_j\) may contain additional edges, triangles, tetrahedra or
higher-dimensional simplices. It need not occupy a larger region of ambient
space.}

For a topological space with a function \(f:X\to\mathbb{R}\), the sublevel set
at threshold \(a\) is
\[
X_a=\{x\in X:f(x)\leq a\}.
\]
Persistent homology often studies the homology of these sublevel sets as
\(a\) increases. In theoretical settings, \(f\) may be given. In data practice,
\(f\) is usually constructed to reveal structure.

All the usual filtrations can be read as sublevel-set filtrations on some
object. Rips, \v{C}ech, lower-star and cubical filtrations differ because they
choose different objects and different filtration functions.
}

\planningbox{Key constructions}{
For a Vietoris--Rips filtration on a metric point cloud,
\[
f(\sigma)=\max_{i,j\in\sigma} d(i,j).
\]
This means a simplex appears when all of its vertices are pairwise close enough.

\qualification{The increasingly dense graph is only the \(1\)-skeleton of the
Vietoris--Rips complex. Every clique is also filled by the corresponding
higher-dimensional simplex. Those filled triangles and higher simplices are
often exactly what kill homology classes.}

For a lower-star filtration based on a vertex function \(g\), define
\[
f(\sigma)=\max_{v\in\sigma} g(v).
\]
A simplex is born when its highest-valued vertex activates. This automatically
satisfies the monotone-on-faces condition. The vertex function does not need to
come from geometry. It could encode intensity, density, activity, behavioural
state or another scalar quantity.
}

\phrase{The filtration function is the modelling decision that decides what it
means for structure to appear.}

\planningbox{Examples}{
The main example is growing balls around points. At small radius there are many
components. At intermediate radius a loop may appear. At large radius that loop
fills.

Make the nested-complex language concrete with three vertices. At one scale,
\(K_i\) may contain the three vertices and three edges, giving an outline
triangle. At a later scale, \(K_j\) contains all of those simplices plus the
two-simplex filling the triangle. Thus \(K_i\subseteq K_j\): the later complex
is made of more simplices, and the new two-simplex kills the earlier
one-dimensional class.

A bad example should also be used: if a triangle appears before one of its
edges, the sublevel object is not a simplicial complex.

Use the Rips--\v{C}ech contrast carefully. A Rips triangle appears when all
pairs are close enough. A \v{C}ech triangle appears when the three balls have a
common intersection. Thus Rips infers higher-order structure from pairwise
conditions, while \v{C}ech requires a genuinely joint geometric condition.
}

\planningbox{DP relevance}{
For a dynamical or collective system, a filtration could be built from distance,
interaction strength, alignment, density or a behavioural score. A lower-star
filtration is useful to keep in mind because vertex values need not be metric.
For example, a vertex value could encode an agent's internal state or local
behavioural feature. This gives a way to study topological structure driven by
behaviour rather than position.
}

\weekheading{4}{Persistent homology}

\planningbox{Story}{
Week 4 closes the first-half definition. Homology is now read across the
filtration using maps induced by inclusion. Persistent homology is not just
``holes at many scales''. It is homology with memory: it records which classes
are born, how long they survive, and when they die.
}

\planningbox{Mathematics to introduce}{
For \(i\leq j\), the inclusion
\[
\iota_{i,j}:K_i\hookrightarrow K_j
\]
induces a linear map on homology:
\[
(\iota_{i,j})_*:H_p(K_i;\mathbb F)\to H_p(K_j;\mathbb F).
\]

The intuition should be made explicit. Every simplex and chain in \(K_i\) is
also present in \(K_j\). A cycle in \(K_i\) can therefore be regarded as the
same chain in \(K_j\), where it is still a cycle. If two cycles differed by a
boundary in \(K_i\), they still differ by that boundary in \(K_j\), so the rule
on homology classes is well defined.

\stickingpoint{The induced map on homology is not necessarily an inclusion.
The larger complex may contain a new higher-dimensional chain that fills an old
cycle. The old chain still exists, but its homology class can map to zero.}

For example, let \(K_i\) be the outline triangle and \(K_j\) the filled
triangle. The loop class is nonzero in \(H_1(K_i)\), but
\[
(\iota_{i,j})_*([z])=0\in H_1(K_j)
\]
because \(z\) has become the boundary of the new two-simplex. This is the
algebraic meaning of death.

In \(H_0\), let two components at scale \(i\) have basis vectors \(e_1,e_2\).
After they merge at scale \(j\), both map to the same component \(e\):
\[
e_1\mapsto e,\qquad e_2\mapsto e.
\]
Hence the difference \(e_1-e_2\) maps to zero. This is the algebra beneath the
merging-components picture and the elder rule.

The persistent homology group from \(i\) to \(j\) is the image
\[
\operatorname{im}\!\bigl(
H_p(K_i;\mathbb F)\xrightarrow{(\iota_{i,j})_*}H_p(K_j;\mathbb F)
\bigr),
\]
and its dimension is the persistent Betti number.

The full \(p\)-dimensional persistence module is the entire collection
\[
V_i=H_p(K_i;\mathbb F),\qquad
\varphi_{i,j}=(\iota_{i,j})_*,
\]
with
\[
\varphi_{i,i}=\operatorname{id},
\qquad
\varphi_{j,k}\circ\varphi_{i,j}=\varphi_{i,k}.
\]

\objectcheck{A filtration is a functor from an ordered parameter set into
simplicial complexes. Homology is a functor from complexes to vector spaces.
Their composition is the persistence module:
\[
(\mathbb R,\leq)\longrightarrow \mathbf{Simp}
\longrightarrow \mathbf{Vect}_{\mathbb F}.
\]
The module is therefore not just the complexes or their inclusions. It is the
homology vector space at every scale together with all compatible linear maps.}

For suitable finite or tame one-parameter persistence modules, there is an
interval decomposition
\[
V\cong \bigoplus_{\ell} I[b_\ell,d_\ell).
\]
A single interval module is
\[
I[b,d)_a=
\begin{cases}
\mathbb F, & b\leq a<d,\\
0, & \text{otherwise},
\end{cases}
\]
with identity maps while the feature is alive. Thus
\[
\text{one interval module}=\text{one bar},
\]
whereas the full persistence module is generally a direct sum of many interval
modules. At any filtration value \(a\), \(\dim V_a\) is the number of bars
crossing \(a\).

\qualification{The barcode does not replace the persistence module by
definition. It records its interval decomposition in the especially
well-behaved one-parameter setting. The finiteness or tameness assumptions are
what make this clean classification possible.}
}

\phrase{A persistence diagram is not just a picture of holes. It is a compressed
record of how homology classes survive through a filtration.}

\planningbox{Stability}{
State the stability result in words: small changes in the filtration function
imply small changes in the persistence diagram. This is the theorem that makes
persistent homology useful for noisy data. But stability is not the same as
interpretability. A stable feature can still be caused by a bad modelling
choice.
}

\planningbox{Examples}{
Build the week around two small calculations before showing a noisy point
cloud.

First, use \(H_0\) for two components merging. Write the source and target vector
spaces and the linear map explicitly, so the audience sees that persistence is
tracking classes rather than merely recounting components at each scale.

Second, use the outline triangle included in the filled triangle. The same
one-cycle is present as a chain in both complexes, but it maps from a nonzero
homology class to zero once the triangular face is available.

Then use the noisy circle as the success case: one long \(H_1\) bar appears
above a cloud of short bars. Use an annulus versus filled disk to show that
persistence is about survival across scale, not presence at a single scale.
Finally, draw several bars and take a vertical slice through them to show that
the vector space at that parameter has one basis direction for each interval
alive there.
}

\planningbox{DP relevance}{
The DP use is initially diagnostic. Rather than asking whether a topological
feature exists at one threshold, ask whether it persists across a meaningful
range of scales. Candidate interpretations might involve fragmentation,
recurrence, circulation, communication pathways or collective organisation. At
this point, the right question is still modest: does static persistence add
information beyond simpler metrics?
}

\section*{Interlude: every barcode has a modelling history}
\addcontentsline{toc}{section}{Interlude: every barcode has a modelling history}

\planningbox{Story}{
The interlude should pause the mathematical build and make the complete modelling
pipeline explicit:
\[
\text{scientific system}
\to \text{observed data}
\to \text{data object}
\to K
\to f
\to K_\bullet
\to \underbrace{H_p(K_\bullet)}_{\text{persistence module}}
\to \text{barcode or diagram}
\to \text{summary, comparison or claim}.
\]
By this point, participants have already encountered most of these choices, but
many will have appeared as though they were simply part of the definition.
The purpose of the interlude is to expose them as modelling decisions. Persistent
homology is not one method applied after the data have been prepared. It is a
family of pipelines in which each decision determines what can later be called
a component, loop, cavity, birth, death or persistent feature.

\stickingpoint{The persistence module is the evolving homology calculation.
The barcode or persistence diagram is the interval-decomposition summary of
that module in the ordinary one-parameter setting. The diagram does not erase
the modelling history that produced the module.}
}

\planningbox{Core message}{
A barcode is never independent of the construction that produced it. The same
observations with a different representation, the same point cloud with a
different metric, the same graph with a different complex, or the same scalar
field with a sublevel rather than superlevel filtration may produce different
diagrams. These are not necessarily competing estimates of one hidden truth.
They may be answers to different scientific questions.
}

\phrase{Before interpreting a bar, reconstruct the sequence of choices that
made that bar possible.}

\planningbox{Design decisions already in play}{
The table should be used actively. For each row, ask participants first to name
the choice made in the examples so far, then to predict what would change under
one alternative. The final column should not be read as a fixed dictionary. It
records the kind of interpretive claim that becomes available, or unavailable,
after that decision.

\begingroup
\small
\renewcommand{\arraystretch}{1.22}
\begin{longtable}{@{}>{\raggedright\arraybackslash}p{0.16\textwidth}
                    >{\raggedright\arraybackslash}p{0.20\textwidth}
                    >{\raggedright\arraybackslash}p{0.25\textwidth}
                    >{\raggedright\arraybackslash}p{0.30\textwidth}@{}}
\toprule
\textbf{Design decision} &
\textbf{Choice used so far} &
\textbf{Some alternatives} &
\textbf{What changes in the interpretation} \\
\midrule
\endfirsthead
\toprule
\textbf{Design decision} &
\textbf{Choice used so far} &
\textbf{Some alternatives} &
\textbf{What changes in the interpretation} \\
\midrule
\endhead
\midrule
\multicolumn{4}{r}{\small\itshape Continued on next page}\\
\endfoot
\bottomrule
\endlastfoot

Scientific object and representation
& A point cloud is treated as the object to be analysed.
& A raw trajectory, delay embedding, interaction graph, scalar field, image,
  time window or state-space sample.
& The resulting topology belongs to the chosen representation. A loop in a
  delay embedding is a statement about recurrent state-space organisation, not
  necessarily a literal loop in physical space. \\

Sampling and preprocessing
& Observed points are used directly and on their given coordinate scales.
& Resampling, denoising, subsampling, standardising coordinates, changing the
  observation window, or retaining temporal order.
& Apparent topology may reflect sampling density, anisotropic units, noise or
  window length. Preprocessing can reveal structure, but can also manufacture
  or remove it. \\

Meaning of closeness
& Euclidean distance on the point cloud.
& Periodic distance, geodesic distance, correlation distance, diffusion
  distance, a learned metric, or a domain-specific dissimilarity.
& Neighbourhoods change, so every later simplex and homology class changes. A
  persistent feature means coherence under the selected notion of closeness,
  not under closeness in general. \\

Combinatorial object $K$
& A Vietoris--Rips complex, or a simplicial complex already supplied.
& \v{C}ech, alpha, witness, cubical, graph, clique or another complex adapted to
  the data type.
& The construction decides which relations are represented explicitly and
  which higher-order relations are inferred. Different complexes can answer
  different questions even when built from the same data. \\

Graph versus clique complex
& Pairwise edges are allowed to generate higher-dimensional Rips simplices.
& Retain only the graph as a one-dimensional complex, or fill every clique to
  form the flag complex.
& A triangle in the graph has an $H_1$ cycle if its interior is not filled. The
  same three pairwise relations have no $H_1$ class in the clique complex.
  Filling cliques asserts that mutual pairwise relation constitutes a
  higher-order unit. \\

Filtration function $f$
& Distance controls a Rips filtration; a vertex value controls a lower-star
  filtration.
& Density, intensity, interaction weight, alignment, energy, activity or
  another scalar attached to vertices, edges or cells.
& The filtration decides what it means for structure to appear. Birth and death
  inherit their scientific meaning from $f$; they do not have a universal
  interpretation. \\

Filtration direction
& Sublevel sets add low-valued parts first.
& Superlevel sets add high-valued parts first, often implemented by filtering
  on $-f$.
& A sublevel filtration detects structures organised around low values; a
  superlevel filtration detects structures organised around high values. Peaks
  and basins exchange roles. \\

Rips scale convention
& A simplex enters when all pairwise distances are at most the reported
  threshold.
& Some accounts grow balls of radius $r$ and use a Rips threshold $2r$; software
  and papers vary in whether the parameter denotes radius or diameter.
& The same filtration can receive birth and death coordinates differing by a
  factor of two. The topology is unchanged, but the numerical scale and any
  physical interpretation of persistence are not. \\

Parameter range and sampling
& The filtration is shown as a finite nested sequence of thresholds.
& A finer or coarser grid, adaptive critical values, a different maximum scale,
  or stopping before all features die.
& Coarse sampling can hide short-lived changes. Truncation can censor deaths and
  make a bar appear essential when it is only unresolved within the chosen
  range. \\

Ties in filtration value
& Simplices with the same filtration value are conceptually added together.
& Software may impose an internal order among tied simplices while respecting
  face-before-coface constraints.
& Arbitrary ordering inside a tie should not be mistaken for scientific time.
  It can affect representatives and create zero-length intermediate events even
  when the persistence module at the stated parameter values is unchanged. \\

Coefficient system
& Homology is computed over $\mathbb{F}_2$.
& Another field $\mathbb{F}_p$, or integer coefficients when torsion and
  orientation matter.
& Over $\mathbb{F}_2$, signs disappear and chains are easy to compute. Changing
  coefficients can alter which classes are visible, especially when torsion is
  present. A Betti number is always relative to the chosen coefficients. \\

Homological dimension $p$
& $H_0$ for components and $H_1$ for loops, with $H_2$ introduced through a
  hollow tetrahedron.
& Compute one degree only, or several degrees with dimension-specific
  constructions and scales.
& Different degrees describe different forms of organisation. A construction
  suitable for components may be computationally or scientifically inadequate
  for cavities. \\

Ordinary or reduced $H_0$
& Ordinary homology leaves one component class alive indefinitely.
& Reduced homology removes the final global component.
& The difference is a convention about the unavoidable connected component,
  not evidence for or against clustering. It changes the essential $H_0$ bar
  and must be checked when comparing software outputs. \\

Object retained after computation
& The persistence module is summarised by its interval decomposition, shown as
  a barcode or persistence diagram.
& Retain the module and its maps, its rank information, selected generators, or
  only a scalar or vector summary.
& For a suitably finite one-parameter module, the barcode classifies the module
  up to isomorphism. It does not automatically retain chosen cycle
  representatives or maps between different modules. The discarded information
  matters later for feature tracking and vineyard modules. \\

Null model and baseline
& The noisy circle and disk are interpreted by visual expectation.
& Shuffled values, random geometric graphs, phase-randomised surrogates,
  resampled point clouds, simpler geometric statistics or domain-specific nulls.
& Persistence alone does not establish that a feature is surprising or useful.
  The null determines what ``more structured than expected'' means, while the
  baseline determines whether topology adds anything beyond a simpler measure.
\\
\end{longtable}
\endgroup

\qualification{Not every row is an independent dial. The data representation
may determine the available complexes; the complex constrains valid filtration
functions; the coefficient field and homological degree affect the algebra; and
the scientific question should constrain all of them. The table is a map of
dependencies, not a menu from which choices should be mixed arbitrarily.}

\objectcheck{At the end of the pipeline, ask which object is actually being
interpreted: a feature of the scientific system, a feature of the observed data,
a class in $H_p(K_a)$, an interval in a persistence module, a point in a diagram,
or a scalar summary of many diagram points. Claims become weaker as more of the
pipeline is forgotten.}
}

\planningbox{Implementation provocations}{
The class should investigate a small selection of the table rows in code. Each
exercise should begin with a prediction, followed by computation, comparison and
an explanation of what changed in the pipeline. The point is not to generate
many diagrams, but to develop the habit of locating the source of a difference.

\begin{enumerate}
    \item \textbf{Graph or filled complex?} Begin with three vertices and all
    three edges. Compute $H_1$ when the object is retained as a graph and when it
    is completed to its clique complex. The data and pairwise relations are
    identical; only the interpretation of a three-way relation changes.

    \item \textbf{Which notion of closeness?} Use one anisotropic point cloud
    and compare Euclidean distance before and after standardising the coordinate
    axes. Alternatively, compare Euclidean and periodic distance for points near
    opposite boundaries of a periodic domain. Identify which apparent
    components or loops were properties of the metric choice.

    \item \textbf{What does birth mean?} Apply lower-star persistence to a
    simple scalar field, then repeat with the negated field. Compare the sublevel
    and superlevel stories and state explicitly whether the features organise
    around low values or high values.

    \item \textbf{What did the software assume?} For a small point cloud, record
    whether the library's Rips parameter is a pairwise-distance threshold or a
    ball radius, whether it returns ordinary or reduced $H_0$, how it represents
    infinite deaths, and how it orders tied simplices. Reconcile any discrepancy
    between two packages before treating it as a mathematical disagreement.

    \item \textbf{Do coefficients matter here?} Compute a supplied triangulation
    with $\mathbb{F}_2$ and with an odd-prime field. A circle should provide the
    control case; a space with torsion, such as a triangulated projective plane,
    shows that coefficient choice can change the answer rather than merely the
    notation.

    \item \textbf{Can topology beat the obvious baseline?} Compare a structured
    example with a suitable surrogate and with one simpler statistic. Require a
    claim of the form: ``under this construction and relative to this null, the
    persistence output reveals \ldots''. If that sentence cannot be completed,
    the pipeline has not yet earned an interpretation.
\end{enumerate}

\stickingpoint{When two outputs differ, the first question should not be ``which
barcode is correct?'' It should be ``which upstream choice changed, and what
question does each pipeline now answer?''}
}

\planningbox{Examples to revisit}{
Return to examples already used in Weeks 1--4 and deliberately alter one choice
at a time:
\begin{itemize}
    \item noisy circle versus noisy disk: change sampling density, metric or
    maximum scale;
    \item outline triangle versus filled triangle: keep the same vertices and
    edges but change whether the two-simplex is included;
    \item merging components: compare ordinary and reduced $H_0$ and identify
    the surviving essential class;
    \item a greyscale field: compare sublevel and superlevel filtrations;
    \item an interaction graph: compare graph-cycle homology with clique-complex
    homology.
\end{itemize}
The repeated examples should make clear that a different answer need not be a
failure of robustness. It may be the mathematically correct consequence of
asking a different question.
}


\weekheading{5}{Static data in practice}

\planningbox{Story}{
Week 5 is the first practical week. The aim is not just to compute persistence,
but to choose a construction appropriate to the data object. The data object
comes first; the complex follows.
}

\planningbox{Method map}{
\begin{itemize}
    \item Point cloud with coordinates: alpha or Rips complexes.
    \item Pairwise distances only: Vietoris--Rips complexes.
    \item Image, volume or gridded scalar field: cubical complexes or lower-star filtrations.
    \item Weighted graph: graph filtration or clique complex.
    \item Large point cloud: sparse Rips or witness constructions.
    \item Labelled or coloured point cloud: chromatic or relative-style constructions.
\end{itemize}
}

\planningbox{Computational cautions}{
Rips complexes grow quickly. For an \(N\)-point cloud, the Vietoris--Rips
filtration contains \(O(N^2)\) edges, \(O(N^3)\) triangles, \(O(N^4)\)
tetrahedra and so on. Even when truncating at dimension 2, the size can scale
like \(O(N^3)\). This is why the choice of maximum dimension, maximum
filtration value and sparsification strategy is not merely technical.
}

\planningbox{Examples}{
Use a noisy circle versus noisy disk as a sanity check. Use an annulus with
uneven sampling to show that persistence is not immune to sampling effects.
Use a greyscale image to introduce lower-star or cubical persistence. Use a
small weighted graph to ask whether pairwise cycles or clique-complex holes are
the object of interest.
}

\planningbox{DP relevance}{
Static persistence is not claiming that a dynamical system is static. It freezes
one representation: a time window, a state cloud, a snapshot, an interaction
graph or a scalar field. This is the lowest-risk baseline. If a single
topological snapshot cannot distinguish regimes or performance-relevant
states, then more sophisticated dynamic topology needs stronger justification.
}

\weekheading{6}{Summaries, comparisons and validation}

\planningbox{Story}{
Once diagrams exist, the next modelling choice is how to use them. A diagram
can be compared directly, summarised by a scalar, converted into a function,
turned into an image, or used as input to statistics and machine learning.
These choices preserve and discard different information.
}

\planningbox{Objects to introduce}{
\begin{itemize}
    \item Bottleneck distance: worst matched feature.
    \item Wasserstein distance: accumulated matching cost.
    \item Betti curves: Betti numbers as functions of filtration scale.
    \item Persistence landscapes: functional summaries in a vector space.
    \item Persistence images: stable vector summaries for machine learning.
    \item Persistence entropy, total persistence, maximum persistence and thresholded counts.
\end{itemize}

\qualification{These objects compare or vectorise persistence diagrams. They do
not, by themselves, recover the maps between persistence modules or establish
the identity of a feature through external time.}
}

\phrase{Stable does not mean automatically interpretable, and vectorised does
not mean scientifically meaningful.}

\planningbox{Examples}{
Use two diagrams with the same longest bar but different clouds of shorter bars
to show why single statistics can be misleading. Use a real-versus-shuffled
comparison to show that the question is not ``is there a barcode?'' but
``does this barcode differ from what this construction would produce anyway?''
}

\planningbox{DP relevance}{
This week asks whether topology adds explanatory or predictive value beyond
simpler baselines. Reasonable first summaries include maximum persistence,
total persistence, Betti curves, persistence entropy, diagram distances to a
baseline and vectorised persistence images or landscapes. Vineyards are only
justified if the identity or trajectory of features matters, not merely their
distribution at each time or parameter value.
}

\weekheading{7}{Time-indexed static topology}

\planningbox{Story}{
Week 7 introduces the honest dynamic baseline. Given a sequence of time-indexed
objects \(X_t\), \(G_t\), \(I_t\), \(W_t\) or \(K_{\bullet,t}\), compute ordinary
persistence independently at each time point or time window. The resulting
diagrams or summaries are then treated as a time series.
}

\planningbox{Mathematics and summaries}{
This is not genuinely dynamic topology in the strongest sense because no
method has identified a specific feature across time. It is repeated static
persistence.

\objectcheck{At each time \(t\), there may be a complete ordinary persistence
module \(V_t\) and its diagram \(D_t\). What is absent is a coherent family of
maps connecting \(V_s\) to \(V_t\). A time series of diagrams is therefore not
yet a vineyard module.}

Useful outputs include
\begin{itemize}
    \item framewise diagrams;
    \item sliding-window diagrams;
    \item Betti curves through scale;
    \item Betti summaries through physical time;
    \item diagram distances to a baseline;
    \item CROCKER-style plots, which project repeated persistence into a rank-function-like summary across scale and time.
\end{itemize}
}

\phrase{A topological summary can change through time without any individual
hole being tracked through time.}

\planningbox{Examples}{
Use moving point clouds, evolving graphs, sliding-window embeddings and
time-indexed scalar fields. The key teaching contrast is between detecting
change in topological summaries and tracking the identity of individual
features.
}

\planningbox{DP relevance}{
This is the first method to try for dynamic systems. It is implementable with
standard tools, interpretable, and often already informative. The DP question
is whether topological summaries over time reveal regime changes, memory,
separability, stability or collective organisation beyond ordinary metrics.
Only if feature identity through time matters should we escalate to more
specialised machinery.
}

\weekheading{8}{Beyond ordinary persistence}

\planningbox{Story}{
Week 8 should be organised by problem type rather than by a list of advanced
methods. The aim is to teach method selection: what breaks in ordinary
persistence, and which tool responds to that break?
}

\planningbox{Case 1: complexes only grow}{
If complexes only grow, ordinary one-parameter persistence is the natural tool.
This is the clean barcode setting. The inclusion maps all point in the same
direction:
\[
K_0\subseteq K_1\subseteq\cdots\subseteq K_n.
\]
}

\planningbox{Case 2: complexes grow and shrink}{
If complexes grow and shrink, zigzag persistence is the natural extension.
Instead of all maps pointing in one direction, the diagram can zigzag:
\[
K_0\to K_1 \leftarrow K_2 \to K_3 \leftarrow \cdots.
\]
Zigzag persistence is designed for insertions and deletions of simplices. It is
mathematically appropriate for non-monotone sequences, but it can be
computationally and interpretively heavy.
}

\phrase{Standard persistence tracks features that survive as a space grows.
Zigzag persistence tracks features that survive as a space changes.}

\planningbox{Case 3: there is more than one genuine parameter}{
If scale and time, scale and density, or another pair of parameters are both
genuine, then multiparameter persistence is the relevant framework. The warning
is essential: the clean barcode decomposition of one-parameter persistence does
not carry over in the same way. Multiparameter persistence may be the right
mathematical object even when a practical analysis later uses lower-dimensional
summaries.

\qualification{The presence of two named variables does not by itself decide
the framework. The question is whether they form one partial order to be
analysed jointly, or whether one parameter indexes a controlled family of
ordinary persistence modules with additional transport data.}
}

\planningbox{Case 4: a fixed complex has a smoothly changing filtration}{
Vineyards apply when there is a fixed simplicial complex \(K\) and a
time-dependent filtration function
\[
f_t:K\to\mathbb{R}.
\]
There are now two distinct parameters:
\[
\underbrace{t}_{\text{external or physical time}},
\qquad
\underbrace{a}_{\text{filtration height or scale}}.
\]
At each fixed time \(t\), the sublevel filtration of \(f_t\) produces an
ordinary persistence module
\[
V_t(a)=H_p(K_{t,a};\mathbb F)
\]
and a persistence diagram \(D_t\). A vineyard tracks how the birth--death
points of \(D_t\) move as \(t\) varies. Each tracked curve is a vine.

The useful teaching picture is a landscape whose height function changes over
time. Ordinary persistence studies flooding one fixed landscape. A vineyard
studies what happens when the landscape itself slowly morphs.

Vineyards assume controlled change. The set of simplices is fixed, filtration
values vary continuously, and changes in the filtration order occur through
adjacent swaps. Between swaps, persistence pairings remain fixed and vines move
linearly or smoothly. At a swap, the reduced boundary matrix is updated locally
rather than recomputed from scratch.

\qualification{A vineyard is a geometric record of moving birth--death points.
It does not automatically retain the full linear maps connecting the
persistence modules at different times.}
}

\phrase{A vineyard is not a stack of unrelated diagrams. It is a one-parameter
family of persistence diagrams realised as curves in birth--death--time space.}

\planningbox{Vineyard modules: what the vineyard omits}{
A vineyard module retains more than the visible trajectories. It consists of
the persistence module \(V_t\) at each external time together with coherent
interleaving maps between the modules at different times. Schematically, the
horizontal direction is ordinary persistence through filtration height, while
the vertical direction transports information through physical time:
\[
\begin{array}{ccccc}
V_s(a)&\longrightarrow&V_s(b)&\longrightarrow&V_s(c)\\
\downarrow&&\downarrow&&\downarrow\\
V_t(a+\varepsilon)&\longrightarrow&
V_t(b+\varepsilon)&\longrightarrow&
V_t(c+\varepsilon).
\end{array}
\]

\stickingpoint{For a single well-behaved one-parameter module, its barcode
classifies the module up to isomorphism. But the source and target barcodes do
not classify a particular map between two modules. Dynamic questions may depend
on those maps.}

A direct sum of vine modules would mean that each visible vine can be treated as
one independent algebraic feature for all time:
\[
I[\gamma_1]\oplus\cdots\oplus I[\gamma_m].
\]
Turner's central point is that a vineyard may visibly contain several vines
while its vineyard module is nevertheless indecomposable. The homology
generators can mix under transport through time.
}

\planningbox{Concrete reading of Turner Figure~1}{
Use Figure~1 of Turner,
\href{https://arxiv.org/abs/2307.06020}{\emph{Representing Vineyard Modules}},
as the concrete example.

The domain is a disc divided into four pieces:
\begin{itemize}
    \item \(A\): the open inner disc;
    \item \(B\): the inner boundary circle;
    \item \(C\): the open surrounding annulus;
    \item \(D\): the outer boundary circle.
\end{itemize}
The time-dependent function assigns
\[
f_t(A)=21-t,\qquad f_t(B)=14-t,\qquad
f_t(C)=15+t,\qquad f_t(D)=t.
\]

\objectcheck{The horizontal axis in the figure is external time \(t\). The
vertical axis is filtration height \(a\). The four black sloping lines are the
values of \(f_t\) on \(A,B,C,D\). At a fixed \(t\), move vertically upward to
watch the sublevel set acquire these four pieces.}

The critical ordering changes when lines cross:
\[
f_3(A)=f_3(C),\qquad f_7(B)=f_7(D).
\]
Hence the dashed vertical lines at \(t=3\) and \(t=7\) divide three regimes.
At every fixed time there are two \(H_1\) bars:
\[
[14-t,21-t)
\qquad\text{and}\qquad
[t,15+t).
\]
The coloured vertical segments drawn at selected times are these bars.
Following each colour through time produces a vine.

The first bar is born when \(B\) enters and dies when \(A\) fills its inner
loop. The second is born when \(D\) enters and dies when \(C\) fills the
annular class.

\stickingpoint{The coloured curves make the example look like two independent
features, but the visible vines do not determine how their homology generators
are transported through time. That missing information is exactly what the
vineyard module retains.}
}

\planningbox{Why the Figure~1 module is indecomposable}{
Before the first critical event, a natural basis is
\[
[B],\qquad[D+B],
\]
whereas after the second critical event the forced basis is
\[
[B],\qquad[D].
\]
Over \(\mathbb F_2\),
\[
[D]=[D+B]+[B].
\]
Thus the change of basis mixes the two visible vines. After coherent
simplification, an off-diagonal term remains in a transport matrix of the form
\[
\begin{pmatrix}
1&1\\
0&1
\end{pmatrix}.
\]

The off-diagonal \(1\) means that transporting the second generator through
time necessarily picks up part of the first. There is no globally coherent
choice of bases that turns the family into two independent \(1\times1\)
problems.

\phrase{Two visible vines do not necessarily mean two independent evolving
homology classes.}

This is the conceptual reason to care about persistence modules rather than
only diagrams: the diagram records where classes are born and die; the
module-level maps record how classes are related and transported.
}

\planningbox{Why vineyards are delicate}{
Smooth motion of points does not necessarily imply smooth behaviour of the
filtration order. In a moving point cloud, many distances can change relative
order at once. Edges, triangles and higher-dimensional simplices can undergo
large global reorderings. This makes it hard to identify the same topological
feature across time. The persistence diagram may be stable in bottleneck
distance, while feature identity is still unstable.

This is the distinction to emphasise:
\[
\text{diagram stability} \neq \text{reliable feature tracking}.
\]
}

\planningbox{Case 5: labels, colours or types matter}{
If points, regions or agents carry labels, colours or types, ordinary
persistence may ignore the structure of interest. Chromatic or relative-style
methods ask how different classes spatially mingle, separate or interact.
Relative homology, image persistence, kernel persistence and cokernel
persistence can be introduced conceptually here as tools for asking what part
of a feature belongs to, is contributed by, or is shared across different
subobjects.
}

\planningbox{DP relevance}{
For swarms or physical reservoirs, vineyards are attractive because they offer
feature identity, continuous trajectories and potential online updates. But
they are also high risk. A literal moving point-cloud Rips filtration may not
satisfy the assumptions needed for vineyard tracking.

Two more plausible strategies are:
\begin{enumerate}
    \item fix an interaction-graph or clique-complex scaffold and let edge or
    simplex weights vary smoothly through time;
    \item fix an ambient mesh and define a smooth density or activity field
    from the swarm or reservoir state.
\end{enumerate}
Both strategies disconnect the fixed simplicial object from the raw moving
agent geometry. The topology then belongs to an interaction structure or
ambient field rather than directly to the moving point cloud.
}

\section*{Closing planning note}
\addcontentsline{toc}{section}{Closing planning note}

The masterclass should not sell advanced TDA as automatically better. The
desired outcome is judgement. Participants should leave able to say:
ordinary persistence is enough here; this is really a zigzag problem; this has
two genuine parameters; this is a labelled point-cloud problem; this is a
possible vineyard problem only if we can define a fixed \(K\) and a controlled
\(f_t\); this question requires the maps between persistence modules rather
than only their diagrams; or topology is not adding enough beyond a simpler
baseline.

That judgement is the foundation needed for the Discovery Project. It is also
the foundation needed for Honours students to use TDA responsibly rather than
as a black-box feature generator.

\end{document}
